\documentclass[10pt,a4paper]{amsart}
\usepackage{amsmath,amssymb,amsfonts}
\usepackage{enumitem}
\usepackage[left=1.5cm, right=1.5cm]{geometry}

\newcommand{\R}{\mathbb{R}}
\renewcommand{\d}{\mathrm{d}}

    \title{Introdução às Equações Diferenciais - Prova 1}

\begin{document}
\maketitle

\noindent {\small {\bf Esta prova é composta por 5 problemas. Leia-os com muita atenção. \\

\noindent Escreva UM problema por folha. É permitido usar mais de uma folha por problema, mas NÃO mais de um problema por folha. \\

\noindent É permitida a utilização de todos os resultados demonstrados em aula, a menos que seja pedido explicitamente que se prove algum destes. Neste caso, é permitida a utilização de todos os resultados vistos em aula ANTERIORES A ESTE. \\

\noindent Tempo total de prova: 3h30min.} }

\vspace{5mm}

\noindent {\bf Problema 1.} Resolva os seguintes problemas de valor inicial. Em cada item, indique explicitamente o maior intervalo em que a solução está definida.
\begin{enumerate}[label=(\alph*)]
    \item ({\bf 5}) $y' - 3y = 6e^{3t}, \qquad y(0) = 1.$
    \item ({\bf 5}) $y' = \dfrac{t}{y}, \qquad y(0) = 2.$
    \item ({\bf 5}) $y' = \dfrac{2t + y}{t}, \qquad y(1) = 1, \qquad t > 0.$ 
    \item ({\bf 10}) 
$\sin (y )\cdot \cos (t) + t \cos (y) \cos (t)\cdot y' = 0, \qquad y\left(\tfrac{\pi}{2}\right) = \tfrac{\pi}{6}.$
\end{enumerate}

\vspace{5mm}

\noindent {\bf Problema 2.} Considere a equação autônoma
\[
    x' = \sin(x).
\]
\begin{enumerate}[label=(\alph*)]
    \item ({\bf 10}) Encontre todas as soluções de equilíbrio e classifique cada uma como estável ou instável.
    \item ({\bf 15}) Suponha que $x(0) = x_0$ com $x_0 \in (0, \pi)$. Prove que a solução $x(t)$ está definida para todo $t \geq 0$.
\end{enumerate}

\vspace{5mm}

\noindent {\bf Problema 3.} Suponha que as funções $\sin t$ e $\sin 2t$ sejam ambas soluções de uma equação diferencial linear homogênea com coeficientes constantes
\[
    x^{(n)} + c_{n-1} x^{(n-1)} + \cdots + c_1 x' + c_0 x = 0,
\]
com $c_0, c_1, \ldots, c_{n-1} \in \R$.
\begin{enumerate}[label=(\alph*)]
    \item ({\bf 15}) Qual é a menor ordem $n$ que essa equação pode ter?
    \item ({\bf 10}) Escreva uma equação diferencial de ordem mínima que tenha ambas as funções como soluções.
\end{enumerate}

\vspace{5mm}

\noindent {\bf Problema 4.} 

\begin{enumerate}
[label=(\alph*)]

\item ({\bf 10}) Prove que toda solução $y$ da equação diferencial
    \[
    y'' + y^3 = 0
    \]
    é limitada. 

    \item ({\bf 15}) Prove que nenhuma solução de 
    \[
    y'' + e^{y} = 0
    \]
    se estende para a reta real toda. 

    \end{enumerate}

\vspace{5mm}

\noindent {\bf Problema 5.} ({\bf 25})  Seja $y:(1,+\infty) \to \mathbb{R}$ uma função continuamente diferenciável tal que 
\[
y'(t) = \frac{t^2 - y(t)^2}{t^2(y(t)^2 + 1)} \quad \text{para todo } t > 1. 
\]
Mostre que $\lim_{t \to \infty} y(t) = + \infty.$
\end{document}